% ============================================================
%  MONOGRAFIA — VERSÃO RESUMO (template UTFPR "RESUMO", v11.0 abr/25)
%  TCC BSI · UTFPR Francisco Beltrão
%  Autor: Felipe Kurt Pohling
%  Orientadora: Prof.ª Dr.ª Mayara Cristina Pereira Yamanoe
%  Fonte do conteúdo: projeto "monografia" (monografia_vFinal)
% ============================================================
\documentclass[12pt,a4paper]{article}

\usepackage{fontspec}
\usepackage{polyglossia}
\setdefaultlanguage{brazil}
\setotherlanguage{english}
\IfFontExistsTF{Arial}{\setmainfont{Arial}}{\setmainfont{TeX Gyre Heros}}

\usepackage[top=3cm,bottom=2cm,left=3cm,right=2cm]{geometry}
\usepackage{setspace}
\usepackage{xcolor}
\usepackage{microtype}
\usepackage[hidelinks]{hyperref}
\hypersetup{
  pdftitle={Comunicação como construção para potencial de pertencimento: o WebNews como artefato sociotécnico},
  pdfauthor={Felipe Kurt Pohling}
}

\setlength{\parindent}{0pt}
\setlength{\parskip}{0pt}
\pagestyle{empty}

% Campos ainda não preenchidos (aparecem em vermelho, como no template)
\newcommand{\pendente}[1]{\textcolor{red}{#1}}

% Dados do trabalho
\newcommand{\autor}{FELIPE KURT POHLING}
\newcommand{\titulo}{COMUNICAÇÃO COMO CONSTRUÇÃO PARA POTENCIAL DE PERTENCIMENTO: O WEBNEWS COMO ARTEFATO SOCIOTÉCNICO}
\newcommand{\tituloingles}{Communication as construction of belonging potential: WebNews as a sociotechnical artifact}
\newcommand{\cidade}{FRANCISCO BELTRÃO}
\newcommand{\ano}{2026}
\newcommand{\natureza}{Trabalho de Conclusão de Curso de Graduação
  apresentado como requisito para obtenção do título de Bacharel em
  Sistemas de Informação, da Universidade Tecnológica Federal do Paraná
  (UTFPR).}

% Bloco de assinatura de membro da banca
\newcommand{\membro}[3]{%
  \begin{center}
    \rule{\linewidth}{0.4pt}\\
    {\small #1\\ #2\\ #3\par}
  \end{center}}

\begin{document}
\frenchspacing

% ============================================================
%  CAPA
% ============================================================
\begin{center}
  \begin{onehalfspace}
  \textbf{UNIVERSIDADE TECNOLÓGICA FEDERAL DO PARANÁ}

  \vspace{3cm}

  \textbf{\autor}

  \vfill

  \textbf{\titulo}

  \vfill
  \vfill

  \textbf{\cidade}\\
  \textbf{\ano}
  \end{onehalfspace}
\end{center}
\clearpage

% ============================================================
%  FOLHA DE ROSTO
% ============================================================
\begin{center}
  \begin{onehalfspace}
  \textbf{\autor}

  \vspace{4cm}

  \textbf{\titulo}

  \vspace{1.5cm}

  \textbf{\tituloingles}
  \end{onehalfspace}
\end{center}

\vspace{1cm}

\hfill
\begin{minipage}{0.5\textwidth}
  \small\singlespacing
  \natureza\\
  Orientadora: Mayara Cristina Pereira Yamanoe.
\end{minipage}

\vfill

\begin{center}
  \begin{onehalfspace}
  \textbf{\cidade}\\
  \textbf{\ano}
  \end{onehalfspace}
\end{center}
\clearpage

% ============================================================
%  FOLHA DE APROVAÇÃO
% ============================================================
\begin{center}
  \begin{onehalfspace}
  \textbf{\autor}

  \vspace{1.5cm}

  \textbf{\titulo}
  \end{onehalfspace}
\end{center}

\vspace{1.5cm}

\hfill
\begin{minipage}{0.5\textwidth}
  \small\singlespacing
  \natureza
\end{minipage}

\vspace{1.5cm}

\begin{center}
  {\small Data de aprovação: 6 de novembro de 2026}
\end{center}

\vfill

\membro{Francisco Antônio Fernandes Reinaldo}
       {Doutorado}
       {Universidade Tecnológica Federal do Paraná}
\vfill
\membro{Flavio de Almeida e Silva}
       {Doutorado}
       {Universidade Tecnológica Federal do Paraná}
\vfill
\membro{Marcos Mincov Tenório}
       {Doutorado}
       {Universidade Tecnológica Federal do Paraná}

\vfill

\begin{center}
  \begin{onehalfspace}
  \textbf{\cidade}\\
  \textbf{\ano}
  \end{onehalfspace}
\end{center}
\clearpage

% ============================================================
%  RESUMO
% ============================================================
\begin{center}
  \begin{onehalfspace}
  \textbf{RESUMO}
  \end{onehalfspace}
\end{center}
\vspace{1cm}

\begin{singlespace}
\raggedright
POHLING, Felipe Kurt. \textbf{Comunicação como construção para potencial
de pertencimento:} o WebNews como artefato sociotécnico. 2026. Trabalho
de Conclusão de Curso (Bacharelado em Sistemas de Informação) --
Universidade Tecnológica Federal do Paraná, Francisco Beltrão, 2026.
\par
\end{singlespace}

\vspace{0.6cm}

\begin{singlespace}
Esta pesquisa investiga a comunicação humana
como mecanismo de construção de potencial de
pertencimento (a disposição pré-vínculo que antecede
a formação de laços sociais genuinamente satisfatórios).
A partir de uma leitura cruzada de dez artigos científicos
em psicologia evolutiva, neurociência, sociologia e
ciências da comunicação, o trabalho demonstra que a
linguagem não é a força primária do processo comunicativo,
mas a ferramenta que o \textit{Mitteilungsbedürfnis}
(a necessidade humana de ser conhecido pelo outro), que invoca
para satisfazer a necessidade fundamental de pertencer;
que as plataformas digitais contemporâneas capitalizam
sistematicamente sobre essa necessidade, produzindo
pertencimento sintético cujos efeitos de
isolamento, polarização e suicídio social são
empiricamente mensuráveis; e que é possível construir
sistemas de informação que operem segundo lógica oposta,
como demonstra o WebNews, plataforma \textit{multi-tenant} de
\textit{newsletters} desenvolvida como artefato empírico desta
pesquisa, cujas decisões de arquitetura, design e modelo
de negócio derivam diretamente dos fundamentos biossociais
da comunicação humana identificados na literatura. A
metodologia combina revisão bibliográfica sistemática,
análise teórica integrativa e \textit{Design Science
Research} (DSR). A contribuição teórica original reside
em três construtos: potencial de pertencimento,
pertencimento sintético e suicídio social,
que emergem da leitura cruzada da literatura e nomeiam
fenômenos empiricamente observáveis mas teoricamente não
consolidados.
\par
\end{singlespace}

\vspace{0.6cm}

\begin{singlespace}
\raggedright
Palavras-chave: pertencimento; comunicação; plataformas digitais;
\textit{newsletter}; sistemas de informação; \textit{Design Science Research}.
\par
\end{singlespace}

\vspace{1.2cm}

(\quad) Não autorizo a disponibilização de endereço de correio eletrônico para contato.

\vspace{0.3cm}

(\,X\,) Autorizo a disponibilização do seguinte correio eletrônico para contato:\\
kvrt.pohling@gmail.com
\clearpage

% ============================================================
%  ABSTRACT
% ============================================================
\begin{english}
\begin{center}
  \begin{onehalfspace}
  \textbf{ABSTRACT}
  \end{onehalfspace}
\end{center}
\vspace{1cm}

\begin{singlespace}
\raggedright
POHLING, Felipe Kurt. \textbf{Communication as construction of belonging
potential:} WebNews as a sociotechnical artifact. 2026. Trabalho de
Conclusão de Curso (Bacharelado em Sistemas de Informação) --
Universidade Tecnológica Federal do Paraná, Francisco Beltrão, 2026.
Título original: Comunicação como construção para potencial de
pertencimento: o WebNews como artefato sociotécnico.
\par
\end{singlespace}

\vspace{0.6cm}

\begin{singlespace}
This research investigates human communication as
a mechanism for constructing belonging potential,
the pre-bond disposition that precedes the formation of
genuinely satisfying social ties. Drawing on a cross-reading
of ten peer-reviewed articles in evolutionary psychology,
neuroscience, sociology, and communication sciences, the
work demonstrates that language is not the primary force
of the communicative process, but the tool invoked by the
\textit{Mitteilungsbedürfnis} (the human need to be known
by others) in order to satisfy the fundamental need to
belong; that contemporary digital platforms systematically
capitalize on this need, producing synthetic
belonging whose effects (isolation, polarization, and
social suicide) are empirically measurable; and that
it is possible to build information systems that operate
according to an opposing logic, as demonstrated by
WebNews, a multi-tenant newsletter platform developed as
the empirical artifact of this research, whose
architectural, design, and business model decisions derive
directly from the biossocial foundations of human
communication identified in the literature. The methodology
combines systematic bibliographic review, integrative
theoretical analysis, and Design Science Research (DSR).
The original theoretical contribution lies in three
constructs: belonging potential,
synthetic belonging, and social suicide,
which emerge from the cross-reading of the literature
and name empirically observable but theoretically
unconsolidated phenomena.
\par
\end{singlespace}

\vspace{0.6cm}

\begin{singlespace}
\raggedright
Keywords: belonging; communication; digital platforms; newsletter;
information systems; Design Science Research.
\par
\end{singlespace}
\end{english}

\end{document}
